\section*{Chapter \ref{ch:nw:the-asia-pacific-map} --- The Asia-Pacific Map}

\begin{solution}{exo:nw:the-asia-pacific-map:1}
Over \qty{6296.4}{\kilo\metre}: \qty{21.00}{\milli\second} in vacuum and \qty{30.71}{\milli\second} in straight fibre, one way.
\end{solution}

\begin{solution}{exo:nw:the-asia-pacific-map:2}
Nine copies later: $9 \times 10 = \qty{90}{\micro\second}$ (105 against \qty{15}{\micro\second} on the first server's port, without jitter).
\end{solution}

\begin{solution}{exo:nw:the-asia-pacific-map:3}
Between Tokyo and Hong Kong the floor is \qty{9.6}{\milli\second}; placing Tokyo within its city moves it by tens of microseconds, under 1\%.
Inside a data centre the competition is over nanoseconds to microseconds, so a building's exact position, and within it the cable
lengths, decide the race.
\end{solution}

\begin{solution}{exo:nw:the-asia-pacific-map:4}
The sender still sends one copy after another at the same cost, so the first and last copies are as far apart as before; what changes is who
receives which copy. A fixed order gives the same member the first copy on every tick; a random one gives each member the first copy on some
ticks and the last on others, so no one has a persistent advantage (rank correlation near zero).
\end{solution}

\begin{solution}{exo:nw:the-asia-pacific-map:5}
Its ports have two members each instead of ten, so even in last place in the centre's order its members receive the tick after
$5 \times 5 + 10$ to $5 \times 5 + 20$, that is 35 to \qty{45}{\micro\second}, ahead of most members of busy servers whatever the shuffle. A
randomiser removes the advantage of logging in first, not the advantage of being on a quieter server: SEBI's order treats the secondary
servers as a separate issue for that reason.
\end{solution}

\begin{solution}{exo:nw:the-asia-pacific-map:6}
Multicast for the strategy: every co-located receiver gets the same copy at the same time, up to the switches' replication and its own cable.
TCP for recovery and for places where multicast does not reach (a disaster-recovery site, a remote office); the FAQs say the secondary data
centre's engine is reached over TCP.
\end{solution}

\begin{solution}{exo:nw:the-asia-pacific-map:7}
\texttt{first\_vs\_median(m)} gives 15.5, 28.5, 50.4 and \qty{100.5}{\micro\second} for $m = 2, 5, 10, 20$, against $c\,m/2 = 10$, 25, 50 and
\qty{100}{\micro\second}. The proposition captures the growth; the small excess at low $m$ comes from the median member's server lying later in
the centre's order, a few copies of \qty{5}{\micro\second}.
\end{solution}

\begin{solution}{exo:nw:the-asia-pacific-map:8}
Equal cables equalise the path from the server to each member, not the time at which the server sends each member its copy. Under sequential
TCP dissemination the tenth member on a port received every tick nine sends after the first, however long its cable.
\end{solution}

\begin{solution}{pb:nw:the-asia-pacific-map:1}
\textbf{Part I.}
\begin{enumerate}
\item First after $5 + 10 = \qty{15}{\micro\second}$, last after $5 + 100 = \qty{105}{\micro\second}$ (plus jitter).
\item The fifth and sixth members receive after 55 and \qty{65}{\micro\second}: \qty{60}{\micro\second} on average.
\item Everything is $3 \times 5 = \qty{15}{\micro\second}$ later: 30 for the first member, 120 for the last.
\item First: the first member on each port of the first server, at \qty{15}{\micro\second}; last: the tenth members of the fourth server, at
  \qty{120}{\micro\second}.
\end{enumerate}
\textbf{Part II.}
\begin{enumerate}[resume]
\item 15.5, 28.5, 50.4 and \qty{100.5}{\micro\second}.
\item Close to $c\,m/2$: 10, 25, 50, 100; the excess is the centre's order.
\item Every time: 100\% in the model, because the order never changes and \qty{50}{\micro\second} dwarfs the reaction noise of
  \qty{5}{\micro\second}.
\item Decisive: races for stale quotes are decided by microseconds or less, so a member \qty{50}{\micro\second} ahead wins nearly every one
  it enters.
\end{enumerate}
\textbf{Part III.}
\begin{enumerate}[resume]
\item The spread stays about \qty{110}{\micro\second}; the rank correlation falls from 1.00 to 0.07; the persistent win rate falls only to 88\%,
  because the quieter secondary server keeps its members ahead.
\item The win rate falls to 51\%: no one is persistently first.
\item The spread falls to the jitter (\qty{8.4}{\micro\second} at the extremes, \qty{0.7}{\micro\second} from first to median) and the win
  rate to 50\%.
\item Multicast: it removes both the persistence and the spread; the other two only redistribute a spread that remains.
\end{enumerate}
\textbf{Part IV.}
\begin{enumerate}[resume]
\item \emph{Named result}: with ten members per port, the first member receives each tick about \qty{50}{\micro\second} before the median member
  under sequential unicast, and about \qty{0.7}{\micro\second} before under multicast; under unicast with a fixed order it wins every race
  against the median member.
\item Facts: the two-level sequential order, three ports per server, login order as the sending order, a lightly loaded secondary server.
  Assumptions: the per-copy costs, the jitter, ten members per port, the reaction-time model.
\item Timestamp the same messages on several connections, ports and servers with synchronised clocks (chapters~4 and~5); compute the spread of
  receive times per message and the rank correlation between messages.
\item Because what matters is the time at which data reaches each member, which depends on the whole chain (sending order, server load,
  cables), not on one piece of it.
\item As a requirement that the communication latency from the dedicated co-location areas to the front-end processors be consistent.
\item Websocket streams from crypto venues (chapter~20), relays in public clouds (chapter~16), and any vendor fan-out over TCP.
\item A shorter path is a fact of geography or engineering, open to anyone who pays for it; a sending order is a choice of the sender,
  which a fair venue must make independent of who the recipient is.
\item Because the feed sent each tick to members one after another in the order they had connected, so the first to connect saw every tick
  first.
\end{enumerate}
\end{solution}

\begin{solution}{iq:nw:the-asia-pacific-map:1}
Unicast sends one copy per recipient over its own connection, one after another; multicast sends one copy to a group address that the network
replicates. In co-location multicast gives every receiver the same copy at the same time and costs the venue one send per message, however
many receivers.
\iqlookfor{Sequential sends and their order; replication in switches; loss and recovery over a separate unicast channel.}
\end{solution}

\begin{solution}{iq:nw:the-asia-pacific-map:2}
Tokyo--Hong Kong about \qty{2900}{\kilo\metre} (\qty{9.6}{\milli\second} in vacuum, \qty{14}{\milli\second} in fibre), Hong Kong--Singapore about
\qty{2600}{\kilo\metre} (\qty{8.6}{\milli\second}), Tokyo--Singapore about \qty{5300}{\kilo\metre} (\qty{17.7}{\milli\second}). Cross-market
strategies there work at millisecond horizons and must choose a home market for each engine.
\iqlookfor{Order of magnitude right; one-way against round trip; the choice of where to put each engine.}
\end{solution}

\begin{solution}{iq:nw:the-asia-pacific-map:3}
Check the capture first (both timestamped by synchronised hardware at the same point); then compare the two connections' servers, ports and
login times, and whether the lag is constant (a sending order), grows with bursts (a slow receiver or a queue) or appears after reconnection;
reconnect in a different order and see whether the lag moves.
\iqlookfor{Measurement before theory; sending order against queueing; an experiment that separates them.}
\end{solution}

\begin{solution}{iq:nw:the-asia-pacific-map:4}
That every recipient of the same service gets each message at the same time up to a spread that does not depend on who it is. Test it by
timestamping the same messages on several connections and checking the spread and whether the order of recipients persists.
\iqlookfor{Both parts: spread and persistence; the venue's published latency reports; equal conditions within a service tier.}
\end{solution}

\begin{solution}{iq:nw:the-asia-pacific-map:5}
Whoever the server writes to first: usually the order of connection or of an internal list, possibly different per update; the last client
waits for all the writes before it, and a slow client can delay the others if the server writes synchronously.
\iqlookfor{Fan-out cost grows with the number of clients; sending order; per-client queues and slow consumers.}
\end{solution}

\begin{solution}{iq:nw:the-asia-pacific-map:6}
Multicast in co-location with equal-length cables from one replication point; one service tier per product with published conditions; a
unicast recovery channel that is not faster than the multicast; per-message timestamps at every member hand-off, and published statistics of
the spread and of the persistence of order.
\iqlookfor{Design and evidence together; the recovery path as a hidden advantage; independent measurement.}
\end{solution}
